% !TeX root = ../papers/nc-p3.tex
\section{整余因子曲線を持つ非恒等平方根の無限小剛性}
\label{sec:conductor-root-tangents}
\begin{theorem}\label{thm:conductor-root-tangents}
定理\ref{thm:integral-conductor-identifiability}の仮定で、二十印を固定した二乗空間の
正規化平方根スキーム $\mathcal F_{X^2}$ は非恒等根の各代表で接空間零である。
従ってそれらは被約孤立点である。
\end{theorem}
\begin{proof}[余因子行列の一次変化]
$k[\epsilon]/(\epsilon^2)$ 上の正規化根を選ぶと、四切断の基底を持ち上げて二つずつの積を
固定された $X^2$ の十個の二次式で一意に表せる。対応行列を $N+\epsilon A$、
余因子の変化を $J$、行列式の変化を $d$ とする。印上で二乗になるので、単元一成分の消去により
$J$ の全成分は二十印の平方イデアルに入る。$C$ の一般点で $N J=0$。
実際 $N\operatorname{adj}N=0$ を $D$ 上で使い、一次変化の恒等式
$NJ+A\operatorname{adj}N=dI$ を制限すればよい。$d=\operatorname{tr}(A\operatorname{adj}N)$
および $\operatorname{adj}N$ は $C$ 上で零である。従って $J$ の像は一般に $N$ の二次元核に入る。

二十次では正規化曲線上で共通固定部を除くと $J$ の成分の束は次数零以下である。
非零なら全体が固定定数行列のスカラー倍となり、その固定像中の非零ベクトルを
$N$ が $C$ 全体で消す。二十印への二次評価の単射性で矛盾する。
二十一次では割線平面六次 $P$ とその印因子 $Z$ を使い、$J$ の成分は $6H-2Z$ の切断で次数は六。
もし $J$ の一般階数が二ならその像二平面は $\ker N$ に一致する。この核は共通二次因子 $q$ の
二つの三次評価の張る割線（非被約なら接線）であり、その直交は $q S_1\subset S_3$。
従ってそのPlücker線束は $2h$ である。実際 $sq,tq$ の外積は二点行列式で割ると $q\otimes q$ になる。
$\bigwedge^2 J$ は対称階数一で、射影的にこのPlücker点の平方である。
左辺は次数十二の束の切断であり、固定部を除いた射の次数は高々十二。
右辺の束は $4h$ で次数二十四である。両写像を正規化上で比較しても束の次数は変わらず、矛盾する。
階数一なら $J$ の与える写像はVeronese像に入る。前節の平面六次の補題により
非定数なら $6H-2Z=h$ で全ての平面一次式を使い、像が平面六次となり矛盾する。
定数像も二十次と同様に排除される。これで $J$ は全て $I_C$ に入る。

$d$ の入力偏導は
$\operatorname{tr}(J\,\partial N+\operatorname{adj}(N)\,\partial A)$ なので $C$ 上で全て零となる。
$C$ の一般点は滑らかで標数零であり、$d$ 自身も零なので $d\in I_C^2$ が一般点で成立する。
前定理の導手交点数から $d$ は元の八次式の定数倍である。ゆえに変形後の行列式は単元倍の同じ超曲面で、
全余因子はその導手イデアルの定数変形に入る。
\end{proof}

\begin{proof}[導手が保たれる変形の正規化作用]
ここからの局所環は双対数での定数拡大を添字 $\epsilon$ で表す。
$M_\epsilon=\operatorname{coker}(N+\epsilon A)$ は基底上平坦である
（自由表示の特殊繊維で最初の写像が単射であることによる）。
$D_\epsilon$ 上では行列分解を持つ極大Cohen--Macaulay加群で、一般には階数一自由である。
非零因子について捩れがなく、全商環中の分数加群とみなせる。Gorenstein環上の極大Cohen--Macaulay双対性から
自然な二重双対写像は同型である。局所的に $A_\epsilon=\mathcal O_{D_\epsilon}$、
$B_\epsilon=(\nu_*\mathcal O_S)_\epsilon$, $I_\epsilon$ を定数導手と書く。
余因子が生成する跡イデアルについて
\[
        M_\epsilon M_\epsilon^\vee\subset I_\epsilon,\qquad
        (B_\epsilon M_\epsilon)M_\epsilon^\vee\subset B_\epsilon I_\epsilon
          \subset A_\epsilon .
\]
全商環中で双対は分数乗法なので、二重双対の等号から $B_\epsilon M_\epsilon\subset M_\epsilon$。
これにより正規化の定数変形上に平坦な層を得る。その特殊繊維は既知の可逆層であり、
局所的に生成元を持ち上げ、平坦性と中山の補題を使えば層自身も可逆である。
$H^1(S,\mathcal O_S)=0$（二次曲面または有理母線束）なので、その直線束の一次変形は自明である。

従って表示加群は定数変形である。最小次数零生成空間と次数二関係を基底変更することで元の対称表示を得る。
実際生成空間は正規化線束の全切断、関係も自由表示の最初の次数で切断を取ればよく、
ともに定数変形である。片側の基底を揃えた対称表示の違いは自己準同型を与え、
その環は $H^0(S_\epsilon,\mathcal O_{S_\epsilon})=k[\epsilon]/(\epsilon^2)$
なので、出力合同と単元スカラーだけである。基底変更を特殊繊維で恒等に取れる
（定数な同型で特殊繊維の選択を合わせる）。各標識点で実際の四切断の平方を読めば同じ四切断に
この共通基底変更とスカラーの平方根を施したものになる。平方からの符号は特殊繊維で固定され、
二が可逆なので無限小には変化しない。従ってGrassmann多様体内の四空間は定数である。
\end{proof}
\begin{proposition}\label{prop:finite-opposite-root-reduced}
階数二以下軌跡が有限な全階数の非恒等根についても、同じ正規化平方根スキームはその根で被約孤立である。
\end{proposition}
\begin{proof}
定理\ref{thm:double-point-opposite-reconstruction}の余因子分解を双対数へ持ち上げる。
変形後の十余因子も固定された $I_{2X}$ に入る。特殊繊維で極小生成系なので
双対数上でも $I_{2X}\otimes k[\epsilon]$ を生成する。
対称行列の同じ複体はこの定数イデアルの自由分解である。
各次数で定数な商まで付けた複体は有限自由加群からなり、特殊繊維で分裂完全、
従って単元小行列で順に消去すれば双対数上でも完全だからである。
最左の行列はしたがって定数変形した元のものに、両側の次数零基底変更を施したものである。
線形化してもこの等式が成立する。十成分の二次座標変換を $g=1+\epsilon d$ と書けば、
$g$ は最左行列の階数三以下の小行列式スキームを保つ。

その特殊繊維の被約台は相手側のVeronese像であり、$d$ はこの台のイデアルも保つ。
実際特性零の微分がイデアルを保てば根基も保つ。$f^n$ がイデアルに入るとき
$n$ 回微分して根基で割ると $n!(df)^n$ が零になるためである。
従って $g$ はVeronese像の定数変形を保つ。二次Veroneseの射影自己同型は双対数上でも
四変数の基底変更から来る。実際像を $\mathbb P^3$ と同一視すると一次束を保ち、
その全切断で基底変更を得る。像は全二次系で埋め込まれているので射影空間の作用もそれで決まる。
従って変形した根行列は合同と単元倍であり、先の証明と同じく正規化根の接方向は零である。
\end{proof}
\begin{corollary}\label{cor:integral-root-collision-exclusion}
この非恒等根へ、固定二乗空間の相対平方根の二つの異なる局所枝が同じ標識付き配置として衝突することはない
（基底は離散付値環、特殊辺は有限被約とする）。
またこの根を対称反復位置に取る反対階数十の片側対称有限窓は全次数正則でない。
\end{corollary}
\begin{proof}
相対の場合は基底の有限拡大後に二乗空間と符号を揃え、
二つのGrassmann切断の最初の差を取ると、特殊繊維で非零接方向となる
（系\ref{cor:opposite-square-root-collision}の証明）。後者は
定理\ref{thm:single-edge-instanton-extension}の非退化接方向の必要性である。
\end{proof}
\path{computations/verify_conductor_root_tangents.py} は階数二での余因子微分と二項割線の全Plücker二次系を記号計算する。
